\documentclass{article}

\usepackage[margin=1in]{geometry}
\usepackage[T1]{fontenc}
\usepackage{lmodern}
\usepackage{amsmath,amssymb,amsthm,mathtools}
\usepackage{array,booktabs,longtable}
\usepackage[hidelinks]{hyperref}
\usepackage[nameinlink,noabbrev]{cleveref}

\providecommand{\lean}[1]{\texttt{#1}}
\providecommand{\leanok}{}
\providecommand{\mathlibok}{}
\providecommand{\uses}[1]{\par\smallskip\noindent\emph{Uses:} #1.\par\smallskip}
\providecommand{\proves}[1]{\par\smallskip\noindent\emph{Proves:} #1.\par\smallskip}
\providecommand{\notready}{\par\smallskip\noindent\emph{Status: placeholder for a later proof phase.}\par\smallskip}

\newtheorem{definition}{Definition}[section]
\newtheorem{lemma}[definition]{Lemma}
\newtheorem{proposition}[definition]{Proposition}
\newtheorem{theorem}[definition]{Theorem}
\newtheorem{corollary}[definition]{Corollary}
\theoremstyle{remark}
\newtheorem{remark}[definition]{Remark}

\newcolumntype{L}[1]{>{\raggedright\arraybackslash}p{#1}}

\newcommand{\C}{\mathbb{C}}
\newcommand{\E}{\mathbb{E}}
\newcommand{\Prb}{\mathbb{P}}
\newcommand{\id}{\mathrm{Id}}
\newcommand{\tr}{\operatorname{tr}}
\newcommand{\Mat}{\operatorname{Mat}}
\newcommand{\code}{\mathcal{C}}
\newcommand{\Hilb}{\mathcal{H}}
\newcommand{\eps}{\varepsilon}

\title{Blueprint for the arXiv:2111.08131 Pilot Formalization}
\author{MIPStarRE}
\date{March 7, 2026}

\begin{document}
\maketitle
\tableofcontents

\section{Scope and conventions}

This document records the initial blueprint for formalizing the tensor-code testing paper arXiv:2111.08131 inside the repository \texttt{MIPStarRE/}. The scope is intentionally narrower than the full paper. The primary goal is a finite-dimensional pilot formalization of the reusable material from Sections~2--3 and the appendix: the matrix-based measurement language, the consistency and closeness calculus, the coding-theoretic interface for tensor codes, the tensor code test, and the expander-graph averaging lemma. Concretely, the pilot works in matrix algebras $\Mat_D(\C)$ with normalized trace, not in general tracial von Neumann algebras.

The secondary goal is to record a theorem and dependency skeleton for Sections~4--6. Those sections depend on the pilot API but also introduce genuinely harder analytic steps: orthogonalization, positive-map duality, the self-improvement mechanism, the pasting construction, and the final induction. These results should therefore appear in Lean first as named placeholders with accurate statements and dependencies, not as immediate proof targets.

The blueprint is organized around two principles. First, reusable mathematics should live in \texttt{MIPStarRE/Quantum/}, \texttt{MIPStarRE/Codes/}, or \texttt{MIPStarRE/Games/}. Second, only the paper-specific induction and the hard Section~4--6 statements should live in \texttt{MIPStarRE/Paper2111/}. The current Lean tree contains only file stubs, so the blueprint intentionally omits \verb|\lean{...}| annotations for now. The dependency tables in Section~\ref{sec:tables} record the intended module home for each node; once declarations exist, each node should receive a corresponding Lean name.

\section{Reusable library targets}

\subsection{Finite-dimensional measurement language}

\begin{definition}[Finite-dimensional tracial ambient space]\label{def:fd-trace}
  For the pilot formalization, the ambient operator algebra is $\Mat_D(\C)$ for a positive integer $D$, equipped with the normalized trace
  \[
    \tau_D(X)=\frac{1}{D} \tr(X).
  \]
  For $X\in\Mat_D(\C)$ we define the trace seminorm
  \[
    \|X\|_\tau = \sqrt{\tau_D(X^*X)}.
  \]
  All later notions of closeness and consistency are specialized to this finite-dimensional setting.
\end{definition}

\begin{proposition}[Finite-dimensional H\"older and triangle bounds]\label{prop:holder}
  \uses{def:fd-trace}
  For all $A,B\in \Mat_D(\C)$ the following hold:
  \[
    |\tau_D(A)| \le \tau_D(|A|),\qquad
    \tau_D(|AB|) \le \|A\|_\tau\,\|B\|_\tau,
  \]
  \[
    \tau_D(|AB|) \le \|A\| \tau_D(|B|),\qquad
    \tau_D(|A+B|) \le \tau_D(|A|)+\tau_D(|B|).
  \]
\end{proposition}
\begin{proof}
  Use polar decomposition in $\Mat_D(\C)$ and the Cauchy--Schwarz inequality for the sesquilinear form $(A,B)\mapsto \tau_D(A^*B)$. The proof is completely finite-dimensional and should be kept separate from any later operator-algebraic generalization.
\end{proof}

\begin{definition}[Submeasurements and processed measurements]\label{def:submeasurement}
  Let $\mathcal A$ be a finite answer set.
  A \emph{submeasurement} on $\mathcal A$ is a family of positive semidefinite matrices $M=\{M_a\}_{a\in\mathcal A}$ such that $\sum_a M_a\le \id$. It is a \emph{measurement} if $\sum_a M_a=\id$, and it is \emph{projective} if each $M_a$ is a projection.

  If $f\colon \mathcal A\to \mathcal B$ is a map to another finite set, the \emph{processed submeasurement} $M_{[f]}$ is defined by
  \[
    M_{[f\mid b]} = \sum_{a:f(a)=b} M_a.
  \]
\end{definition}

\begin{definition}[Consistency and closeness]\label{def:consistency}
  Let $\mathcal X$ be a finite question set with distribution $\mu$, let $\mathcal A$ be a finite answer set, and let $M^x=\{M^x_a\}_a$ and $N^x=\{N^x_a\}_a$ be submeasurements indexed by $x\in\mathcal X$.

  We say that $M$ and $N$ are \emph{$\delta$-consistent}, and write $M^x_a\simeq_\delta N^x_a$, if
  \[
    \E_{x\sim\mu}\sum_{a\ne b}\tau_D(M^x_aN^x_b)\le \delta.
  \]
  We say that $M$ and $N$ are \emph{$\delta$-close}, and write $M^x_a\approx_\delta N^x_a$, if
  \[
    \left(\E_{x\sim\mu}\sum_a \|M^x_a-N^x_a\|_\tau^2\right)^{1/2}\le \delta.
  \]
\end{definition}

\begin{proposition}[Data processing for consistency]\label{lem:data-processing}
  \uses{def:submeasurement, def:consistency}
  If $M^x_a\simeq_\delta N^x_a$ and $f\colon \mathcal A\to \mathcal B$, then the processed families satisfy
  \[
    M^x_{[f\mid b]}\simeq_\delta N^x_{[f\mid b]}.
  \]
\end{proposition}
\begin{proof}
  Expand the processed outcome classes. Merging answers cannot create new off-diagonal mass, so the inconsistency only decreases.
\end{proof}

\begin{proposition}[Consistency to closeness]\label{lem:consistency-consequences}
  \uses{def:consistency}
  If $M^x$ and $N^x$ are measurements and $M^x_a\simeq_\delta N^x_a$, then
  \[
    M^x_a\approx_{\sqrt{2\delta}} N^x_a.
  \]
\end{proposition}
\begin{proof}
  Expand $\sum_a\|M^x_a-N^x_a\|_\tau^2$, use $\sum_a M^x_a=\sum_a N^x_a=\id$, and rewrite the diagonal overlap in terms of inconsistency.
\end{proof}

\begin{proposition}[Closeness to consistency]\label{lem:closeness-to-consistency}
  \uses{def:consistency}
  Let $M^x$ be a projective submeasurement and let $N^x$ be a submeasurement with the same answer set. If $M^x_a\approx_\delta N^x_a$, then
  \[
    M^x_a\simeq_\delta N^x_a.
  \]
\end{proposition}
\begin{proof}
  Rewrite the off-diagonal inconsistency as $\sum_a \tau(M^x_a(\id-N^x_a))$ and bound this by Cauchy--Schwarz together with projectivity of $M^x$.
\end{proof}

\begin{proposition}[Closeness preserves averaged overlaps]\label{lem:closeness-to-close-ips}
  \uses{def:consistency}
  Let $A^x_a$ and $B^x_a$ be operator families such that $A^x_a\approx_\delta B^x_a$. If $M^x_a$ is another family satisfying $\sum_a M^x_a(M^x_a)^*\le \id$ for every $x$, then
  \[
    \E_x\sum_a \tau_D(M^x_aA^x_a)
    \approx_\delta
    \E_x\sum_a \tau_D(M^x_aB^x_a).
  \]
\end{proposition}
\begin{proof}
  Apply Cauchy--Schwarz to the averaged pairing between $M^x_a$ and $A^x_a-B^x_a$.
\end{proof}

\begin{proposition}[Left and right multiplication preserve closeness]\label{lem:add-a-proj}
  \uses{def:consistency}
  Let $M_a$ and $N_a$ be operator families with $M_a\approx_\delta N_a$. Let $R_b$ be operators satisfying either $\sum_b R_b^*R_b\le \id$ or $\sum_b R_bR_b^*\le \id$. Then left or right multiplication by the $R_b$ preserves the same closeness bound.
\end{proposition}
\begin{proof}
  Expand the squared trace norm after multiplication and use positivity together with the assumption that the $R_b$ contract on average.
\end{proof}

\begin{proposition}[Transfer consistency using closeness]\label{lem:transfer-cons}
  \uses{def:consistency, lem:closeness-to-close-ips}
  Let $A^x$ and $B^x$ be measurements and let $C^x$ be a submeasurement with the same answer set. If $A^x_a\simeq_\delta C^x_a$ and $A^x_a\approx_\eps B^x_a$, then
  \[
    B^x_a\simeq_{\delta+\eps} C^x_a.
  \]
\end{proposition}
\begin{proof}
  Rewrite inconsistency against a submeasurement as total mass minus diagonal overlap. The total mass is unchanged, and the diagonal overlap changes by at most $\eps$ by \Cref{lem:closeness-to-close-ips}.
\end{proof}

\begin{proposition}[Consistency with a measurement controls a submeasurement]\label{lem:cons-sub-meas}
  \uses{def:consistency}
  Let $A^x$ be a submeasurement and $B^x$ a measurement such that $A^x_a\simeq_\gamma B^x_a$. Then
  \[
    A^x_a\approx_{\sqrt\gamma} A^x_aB^x_a
    \qquad\text{and}\qquad
    A^x_a\approx_{2\sqrt\gamma} A^xB^x_a,
  \]
  where $A^x=\sum_a A^x_a$.
\end{proposition}
\begin{proof}
  The inconsistency bound controls the mass of $A^x_a(\id-B^x_a)$ and of $(A^x-A^x_a)B^x_a$. Expanding the squared trace norm gives the stated closeness bounds.
\end{proof}

\begin{proposition}[Switcheroo for approximately commuting families]\label{lem:switcheroo}
  \uses{lem:add-a-proj}
  Let $A^x=\{A^x_a\}$ be submeasurements such that, on average over $x,y$,
  \[
    A^x_aA^y_b\approx_\eps A^y_bA^x_a.
  \]
  Then for every $k\ge 1$, if $P^s_{\vec a}$ denotes a length-$k$ product of operators drawn from the families $A^{s_i}$, one has
  \[
    P^s_{\vec a}A^x_b\approx_{k\eps} A^x_bP^s_{\vec a}
  \]
  on average over $x$ and the sequence $s$.
\end{proposition}
\begin{proof}
  Induct on the product length. The inductive step inserts one extra factor and uses the triangle inequality together with \Cref{lem:add-a-proj}.
\end{proof}

\subsection{Linear and tensor codes}

\begin{definition}[Linear code]\label{def:code}
  Let $\Sigma$ be a finite field. A linear $[n,k,d]_\Sigma$ code is a $k$-dimensional linear subspace $\code\subseteq (\Sigma^n)$ such that any two distinct codewords differ in at least $d$ coordinates.
\end{definition}

\begin{proposition}[Uniqueness from distance]\label{prop:distance0}
  \uses{def:code}
  Let $\code$ be an $[n,k,d]_\Sigma$ code and put $t=n-d+1$. If $i_1,\dots,i_t\in [n]$ are distinct, then for any values $a_1,\dots,a_t\in\Sigma$ there is at most one codeword $c\in\code$ such that $c(i_j)=a_j$ for all $j$.
\end{proposition}
\begin{proof}
  Two distinct codewords agreeing on those $t$ coordinates would disagree on at most $d-1$ coordinates, contradicting the distance bound.
\end{proof}

\begin{definition}[Interpolable code]\label{def:interpolable}
  \uses{def:code, prop:distance0}
  Let $\code$ be an $[n,k,d]_\Sigma$ code and put $t=n-d+1$. We say that $\code$ is \emph{interpolable} if for every choice of distinct coordinates $i_1,\dots,i_t\in [n]$ there is a linear map
  \[
    \operatorname{interp}_{i_1,\dots,i_t}\colon \Sigma^t\to \code
  \]
  sending $(a_1,\dots,a_t)$ to the unique codeword whose values at the chosen coordinates are $a_1,\dots,a_t$.
\end{definition}

\begin{definition}[Axis-parallel line]\label{def:axis-line}
  For $m\ge 1$, an \emph{axis-parallel line} in $[n]^m$ is a subset obtained by fixing all but one coordinate and letting the remaining coordinate vary over $[n]$.
\end{definition}

\begin{definition}[Tensor code]\label{def:tensor-code}
  \uses{def:axis-line, def:code}
  Let $\code$ be an $[n,k,d]_\Sigma$ code and let $m\ge 1$. The tensor code $\code^{\otimes m}$ is the set of functions $c\colon [n]^m\to \Sigma$ such that the restriction of $c$ to every axis-parallel line is a codeword of $\code$.
\end{definition}

\begin{proposition}[Distance of the tensor code]\label{prop:distance}
  \uses{def:tensor-code, prop:distance0}
  Let $t=n-d+1$ and set $\gamma_m=1-d^m/n^m$. For distinct $c,c'\in \code^{\otimes m}$,
  \[
    \Prb_{u\in [n]^m}[c(u)=c'(u)]\le \gamma_m\le \frac{mt}{n}.
  \]
\end{proposition}
\begin{proof}
  The formal proof should establish that the distance of $\code^{\otimes m}$ is at least $d^m$, for example by induction on $m$ or by a generator-matrix argument. The paper's informal sentence that tensor-code codewords are ``exactly the tensor products of codewords'' should not be copied literally into Lean. The estimate $\gamma_m\le mt/n$ then follows from Bernoulli's inequality.
\end{proof}

\begin{proposition}[Tuple-to-code correspondence]\label{prop:tuple-to-code-correspondence}
  \uses{def:tensor-code, prop:distance0}
  Let $x_1,\dots,x_t\in [n]$ with $t\ge n-d+1$. A tuple $(g_1,\dots,g_t)$ of codewords in $\code^{\otimes m}$ comes from at most one codeword $h\in \code^{\otimes (m+1)}$ with $h|_{x_j}=g_j$ for every $j$.
\end{proposition}
\begin{proof}
  If two such $(m+1)$-dimensional codewords existed, then on some line in the last coordinate they would give distinct base-code codewords agreeing on at least $t$ points, contradicting \Cref{prop:distance0}.
\end{proof}

\begin{proposition}[Interpolation of tuples]\label{prop:interpolate-tuple}
  \uses{def:interpolable, prop:tuple-to-code-correspondence}
  If the base code $\code$ is interpolable and $t=n-d+1$, then every tuple $(g_1,\dots,g_t)\in (\code^{\otimes m})^t$ extends to a codeword in $\code^{\otimes (m+1)}$.
\end{proposition}
\begin{proof}
  Apply the interpolation map pointwise in the last coordinate. Linearity of the interpolation map shows that the slices in the first $m$ coordinates remain codewords of $\code^{\otimes m}$.
\end{proof}

\subsection{Expander averaging support}

\begin{proposition}[Laplacian as an edge average]\label{prop:laplacian-edge-form}
  Let $G$ be a finite graph on $n$ vertices with a probability distribution on edges whose vertex marginals are uniform. If
  \[
    K = \E_{(u,v)\sim G} |u\rangle\langle v|,
    \qquad
    L = \frac{1}{n}\id - K,
  \]
  then
  \[
    L = \frac12\,\E_{(u,v)\sim G} (|u\rangle-|v\rangle)(\langle u|-\langle v|).
  \]
\end{proposition}
\begin{proof}
  Expand the edge average and use the uniform vertex marginals to identify the diagonal part with $n^{-1}\id$.
\end{proof}

\begin{lemma}[Local-to-global expander inequality]\label{lem:local-to-global-expander}
  \uses{prop:laplacian-edge-form}
  Let $G$ be as above and let $\lambda_2$ be the second eigenvalue of its normalized Laplacian. For any family of matrices $A^u\in \Mat_D(\C)$ and any positive linear functional $\rho$ on $\Mat_D(\C)$,
  \[
    \E_{u,v\sim [n]} \rho((A^u-A^v)^*(A^u-A^v))
    \le
    \frac{1}{n\lambda_2}
    \E_{(u,v)\sim G} \rho((A^u-A^v)^*(A^u-A^v)).
  \]
\end{lemma}
\begin{proof}
  Form the block operator $V=\sum_u |u\rangle\otimes A^u$, compare $V^*(L\otimes \id)V$ with the projector onto the orthogonal complement of the constant vector, and then apply the positive functional $\rho$.
\end{proof}

\section{Paper-specific targets for the tensor code test}

\begin{definition}[Subcube]\label{def:subcube}
  A \emph{subcube} of $[n]^m$ is obtained by fixing a suffix of the coordinates and allowing the remaining prefix to vary freely. In particular, the whole cube and every axis-parallel line are special cases.
\end{definition}

\begin{definition}[Augmented tensor code test]\label{def:tensor-code-test}
  \uses{def:axis-line, def:subcube, def:tensor-code}
  The augmented tensor code test for $\code^{\otimes m}$ is the synchronous two-question test with two equally weighted branches.

  In the first branch, the referee samples a random point $u\in [n]^m$ and a random axis-parallel line $\ell$ through $u$, asks one prover for a codeword on $\ell$ and the other prover for the symbol at $u$, and checks agreement at $u$.

  In the second branch, the referee samples a random subcube $H$, random points $u,v\in H$, asks for a joint answer on $(u,v)$ and a single answer on $u$, and checks consistency. This branch is designed to force approximate commutation of point measurements inside random subcubes.
\end{definition}

\begin{definition}[Finite-dimensional tracial strategy]\label{def:tracial-strat}
  \uses{def:tensor-code-test}
  A finite-dimensional tracial strategy for the tensor code test consists of a matrix algebra $\Mat_D(\C)$ with normalized trace, together with a projective point measurement $A^u$ for each $u\in [n]^m$, a projective line measurement $B^\ell$ for each axis-parallel line $\ell$, and a projective pair measurement $P^{u,v}$ for each admissible pair question $(u,v)$. The acceptance probability is computed using the normalized trace pairing, exactly as in the synchronous finite-dimensional interpretation of the paper.
\end{definition}

\begin{definition}[$(\eps,\delta)$-good strategy]\label{def:tracial-good}
  \uses{def:tracial-strat, def:consistency}
  A finite-dimensional tracial strategy is \emph{$(\eps,\delta)$-good} if the processed line measurements are $\eps$-consistent with the point measurements on average over a random line and a random point on that line, and the processed pair measurements are $\delta$-consistent with the corresponding point measurements on average over the random subcube branch.
\end{definition}

\begin{proposition}[Winning probability versus goodness]\label{prop:test-goodness}
  \uses{def:tensor-code-test, def:tracial-good}
  If a strategy is $(\eps,\delta)$-good, then it wins the augmented tensor code test with probability at least $1-\tfrac12(\eps+\delta)$. Conversely, if it wins with probability $1-\eta$, then after distributing the two branches one obtains an $(O(\eta),O(\eta))$-good strategy.
\end{proposition}
\begin{proof}
  This is bookkeeping: unfold the two branches of the referee distribution and rewrite each acceptance condition as a consistency statement.
\end{proof}

\section{Section~4--6 theorem skeleton}

The nodes in this section are intentionally coarser than the full paper proofs. They expose the later work packages without pretending that the current pilot will prove them immediately. Nodes marked \notready should initially be introduced in Lean as named placeholders in \texttt{MIPStarRE/Paper2111/Skeleton.lean}. Nodes left unmarked are reasonable first post-pilot proof targets.

\subsection{Self-improvement chain}

\begin{lemma}[Orthogonalization]\label{lem:projectivization}
  \uses{def:submeasurement, def:fd-trace}
  \notready
  If a submeasurement $A=\{A_a\}$ is almost projective in the sense that
  \[
    \sum_a \tau_D(A_a(\id-A_a))
  \]
  is small, then there exists a projective submeasurement $P=\{P_a\}$ with $A_a\approx P_a$ and an explicit square-root error bound.
\end{lemma}
\begin{proof}
  In finite dimensions this should be proved directly by polar decomposition or matrix orthogonalization. Until that argument is written, keep this as a placeholder rather than importing the full infinite-dimensional theorem.
\end{proof}

\begin{lemma}[Dominating positive functional]\label{lem:duality}
  \uses{def:fd-trace}
  \notready
  Given finitely many positive linear functionals $\varphi_1,\dots,\varphi_k$ on $\Mat_D(\C)$, there is a least positive linear functional $\psi$ dominating all of them, together with a measurement $\{T_i\}$ such that
  \[
    \psi(X)=\sum_i \varphi_i(T_iX)
  \]
  for every $X$.
\end{lemma}
\begin{proof}
  In the finite-dimensional pilot, this should eventually be phrased and proved as an explicit semidefinite-program duality statement. For now it should remain a paper-specific placeholder.
\end{proof}

\begin{lemma}[Local variance along a random line]\label{lem:local-variance}
  \uses{def:tracial-good, lem:consistency-consequences, prop:distance}
  For an $(\eps,\delta)$-good strategy and any measurement $\{T_g\}$ on $\code^{\otimes m}$, the sandwiched operators $A^x_{g(x)}T_g^{1/2}$ and $A^y_{g(y)}T_g^{1/2}$ are close on average when $x$ and $y$ lie on a random axis-parallel line.
\end{lemma}
\begin{proof}
  Compare the point measurements with the corresponding processed line measurement, then use the code-distance bound from \Cref{prop:distance} to replace agreement at one point on the line by agreement with the unique line codeword.
\end{proof}

\begin{lemma}[Global variance]\label{lem:variance}
  \uses{lem:local-variance, lem:local-to-global-expander}
  For an $(\eps,\delta)$-good strategy and any measurement $\{T_g\}$ on $\code^{\otimes m}$, the same sandwiched operators are close on average over two uniformly random points of $[n]^m$.
\end{lemma}
\begin{proof}
  Apply \Cref{lem:local-to-global-expander} to the graph on $[n]^m$ connecting points that differ in at most one coordinate. The local estimate from \Cref{lem:local-variance} supplies the edge bound.
\end{proof}

\begin{lemma}[Self-improvement]\label{lem:self-improvement}
  \uses{lem:projectivization, lem:duality, lem:variance, lem:transfer-cons, lem:cons-sub-meas}
  \notready
  Suppose a strategy for the tensor code test already has a global measurement $G$ on $\code^{\otimes m}$ that is reasonably consistent with the point measurements. Then one can construct a projective submeasurement $H$ on $\code^{\otimes m}$ that is nearly complete, still consistent with the point measurements, and whose incomplete part is controlled by a dominating positive functional.
\end{lemma}
\begin{proof}
  Average the point measurements against the current global measurement, choose the dominating functional from \Cref{lem:duality}, define the provisional submeasurement by sandwiching with the point projectors, and then orthogonalize using \Cref{lem:projectivization}. The variance estimate is what keeps the error from blowing up.
\end{proof}

\subsection{Pasting chain}

\begin{lemma}[Commutativity of point measurements]\label{lem:a-comm}
  \uses{def:tracial-good, lem:consistency-consequences, lem:add-a-proj}
  Passing the subcube branch of the tensor code test with high probability implies that two independently sampled point operators $A^{u,x}_a$ and $A^{v,y}_b$ approximately commute on average.
\end{lemma}
\begin{proof}
  The pair measurement witnesses a common refinement of the two point measurements. Convert consistency with that refinement into closeness and then move the approximations through products with \Cref{lem:add-a-proj}.
\end{proof}

\begin{lemma}[Evaluated subspace measurements commute]\label{lem:g-comm-after-eval}
  \uses{lem:a-comm, lem:cons-sub-meas, lem:closeness-to-close-ips}
  \notready
  If the family of subspace measurements $G^x$ is consistent with the point measurements, then the processed families $G^x_{[g\mapsto g(u)\mid a]}$ approximately commute on average over two random points.
\end{lemma}
\begin{proof}
  Replace each evaluated $G$ by the corresponding point measurement using consistency, commute the point measurements by \Cref{lem:a-comm}, and then transfer the result back to the evaluated $G$ operators.
\end{proof}

\begin{lemma}[Subspace measurements commute]\label{lem:g-comm}
  \uses{lem:g-comm-after-eval, prop:distance}
  \notready
  The subspace measurements $G^x_g$ and $G^y_h$ approximately commute on average over independent $x,y\in [n]$.
\end{lemma}
\begin{proof}
  Average the evaluated commutation relation over random evaluation points and use the tensor-code distance bound from \Cref{prop:distance} to control the contribution from distinct codewords that accidentally agree at a sampled point.
\end{proof}

\begin{corollary}[Commutativity of the completed measurement]\label{cor:G-hat-facts}
  \uses{lem:g-comm}
  \notready
  After adjoining the incomplete outcome $\bot$, the completed measurements $\widehat G^x$ also approximately commute.
\end{corollary}
\begin{proof}
  Expand the commutator between complete and incomplete parts and reduce each term to the commutation statements already available for $G^x_g$ and $G^x=\sum_g G^x_g$.
\end{proof}

\begin{lemma}[Consistency of the pasted measurement with line answers]\label{lem:h-b-consistency}
  \uses{cor:G-hat-facts, prop:interpolate-tuple, lem:data-processing, lem:transfer-cons}
  \notready
  The sandwiched-and-interpolated measurement $H$ produced by the pasting construction is approximately consistent with the line measurement $B$.
\end{lemma}
\begin{proof}
  Commute the $\widehat G$ factors until one can expose a single line answer, then use interpolation uniqueness from \Cref{prop:interpolate-tuple} to identify inconsistent tuples and bound them by the line-versus-point consistency error.
\end{proof}

\begin{corollary}[Consistency of the pasted measurement with point answers]\label{cor:ha-cons}
  \uses{lem:h-b-consistency, lem:data-processing, lem:transfer-cons}
  \notready
  The pasted measurement $H$ is approximately consistent with the point measurements.
\end{corollary}
\begin{proof}
  Data-process the line answer down to its value at a random point and then transfer consistency from the line measurement to the point measurement.
\end{proof}

\begin{lemma}[All outcomes versus globally consistent outcomes]\label{lem:over-all-outcomes}
  \uses{prop:interpolate-tuple, prop:distance}
  \notready
  The completeness of the pasted measurement $H$ is close to the total mass of all sandwiched outcome tuples whose type contains at least $t$ successful slices, even if one temporarily drops the requirement that those tuples already come from a global codeword.
\end{lemma}
\begin{proof}
  Use a line measurement as a witness and show that a tuple that is locally consistent too often but not globally consistent would violate the tensor-code distance bound.
\end{proof}

\begin{lemma}[Collapse to a binomial operator polynomial]\label{lem:from-H-to-G}
  \uses{cor:G-hat-facts}
  \notready
  The total mass of sandwiched outcome tuples of type at least $t$ is close to the operator polynomial
  \[
    \sum_{i=t}^k \binom{k}{i} G^i(\id-G)^{k-i},
  \]
  where $G$ denotes the average completeness operator of the subspace measurements.
\end{lemma}
\begin{proof}
  Commute the $\widehat G$ factors until the sandwiched product becomes symmetric, then peel off one coordinate at a time and use the recurrence for the polynomial in $G$ and $\id-G$.
\end{proof}

\begin{lemma}[Matrix Chernoff bound for a Bernoulli polynomial]\label{lem:chernoff-bernoulli-matrix}
  \uses{def:fd-trace}
  \notready
  Let
  \[
    F(x)=\sum_{i=t}^k \binom{k}{i}x^i(1-x)^{k-i}.
  \]
  If $0\le X\le \id$ is self-adjoint and $\tau_D(X)\ge 1-\kappa$, then $\tau_D(F(X))$ is bounded below by $1-\kappa/(1-\theta)$ minus an exponentially small Chernoff term, for every admissible choice of $\theta$.
\end{lemma}
\begin{proof}
  In finite dimensions, diagonalize $X$ and reduce to the scalar additive Chernoff bound on each eigenvalue. This is a clean matrix theorem, but it should remain deferred until the rest of the pilot API is in place.
\end{proof}

\begin{corollary}[Completeness of the pasted measurement]\label{cor:ld-pasting-N-completeness}
  \uses{lem:over-all-outcomes, lem:from-H-to-G, lem:chernoff-bernoulli-matrix}
  \notready
  For a suitable repetition parameter $k\ge 12mt$, the pasted measurement $H$ has completeness at least
  \[
    1-\kappa\left(1+\frac{1}{3m}\right)-\text{poly}(m,t,k) \text{poly}(\eps,\delta,\zeta,n^{-1})-e^{-\Omega(k/m^2)}.
  \]
\end{corollary}
\begin{proof}
  Combine \Cref{lem:over-all-outcomes} and \Cref{lem:from-H-to-G} to reduce the completeness of $H$ to the matrix polynomial from \Cref{lem:chernoff-bernoulli-matrix}, then apply the operator Chernoff estimate.
\end{proof}

\begin{lemma}[Pasting]\label{lem:pasting}
  \uses{cor:ha-cons, cor:ld-pasting-N-completeness}
  \notready
  If each parallel $m$-dimensional slice carries a sufficiently complete and sufficiently point-consistent submeasurement $G^x$, then one can paste these slice measurements into a measurement $H$ on $\code^{\otimes (m+1)}$ that is again consistent with the point measurements, with controlled loss in the error parameter.
\end{lemma}
\begin{proof}
  The consistency part comes from \Cref{cor:ha-cons}; the completeness part comes from \Cref{cor:ld-pasting-N-completeness}. The construction is the sandwiched interpolation procedure from Section~6.
\end{proof}

\subsection{Main induction}

\begin{lemma}[Induction step]\label{lem:induction}
  \uses{lem:self-improvement, lem:pasting}
  \notready
  For every $m\ge 1$, every $(\eps,\delta)$-good strategy for the tensor code test on $\code^{\otimes m}$ admits a global measurement $G$ on $\code^{\otimes m}$ that is approximately consistent with the point measurements, with an explicit polynomial error bound.
\end{lemma}
\begin{proof}
  The induction constructs slice measurements in dimension $m$, pastes them into a measurement in dimension $m+1$ by \Cref{lem:pasting}, and then applies \Cref{lem:self-improvement} to recover projectivity and bounded incomplete mass.
\end{proof}

\begin{theorem}[Finite-dimensional quantum soundness skeleton]\label{thm:main}
  \uses{lem:induction, lem:self-improvement}
  \notready
  For a finite-dimensional tracial strategy that wins the augmented tensor code test for $\code^{\otimes m}$ with high probability, there exists a projective measurement on $\code^{\otimes m}$ whose outcome is approximately consistent with the prover's point measurement at a uniformly random point.
\end{theorem}
\begin{proof}
  Apply \Cref{lem:induction} to obtain a global measurement and then apply \Cref{lem:self-improvement} one final time to make it projective.
\end{proof}

\begin{definition}[Good two-prover strategy]\label{def:good-2}
  \uses{def:tensor-code-test, def:tracial-good}
  A good two-prover strategy for the later extension consists of separate measurement families for the two provers together with a high success probability in the line branch, the subcube branch, and an added synchronicity branch that forces the two provers to answer identically when they receive the same question.
\end{definition}

\begin{theorem}[Two-prover extension skeleton]\label{thm:main-bipartite}
  \uses{thm:main, def:good-2}
  \notready
  A near-optimal two-prover strategy for the synchronized tensor code test is close to a pair of global codeword measurements, one on each prover side, that agree with each other and with the original point and line measurements.
\end{theorem}
\begin{proof}
  Reduce the two-prover strategy to the synchronous finite-dimensional setting and then apply \Cref{thm:main}. The reduction should remain a placeholder until the finite-dimensional pilot is stable.
\end{proof}

\section{Dependency tables and work phases}\label{sec:tables}

\subsection{Pilot nodes}

\small
\begin{longtable}{L{0.16\textwidth}L{0.27\textwidth}L{0.33\textwidth}L{0.14\textwidth}}
\toprule
Node & Intended Lean home & Direct dependencies & Initial policy \\
\midrule
\endfirsthead
\toprule
Node & Intended Lean home & Direct dependencies & Initial policy \\
\midrule
\endhead
\texttt{def:fd-trace} & \texttt{MIPStarRE/Quantum/Measurement.lean} & --- & prove now \\
\texttt{prop:holder} & \texttt{MIPStarRE/Quantum/Measurement.lean} & \texttt{def:fd-trace} & prove now \\
\texttt{def:submeasurement} & \texttt{MIPStarRE/Quantum/Measurement.lean} & \texttt{def:fd-trace} & prove now \\
\texttt{def:consistency} & \texttt{MIPStarRE/Quantum/Measurement.lean} & \texttt{def:submeasurement} & prove now \\
\texttt{lem:data-processing} & \texttt{MIPStarRE/Quantum/Measurement.lean} & \texttt{def:submeasurement}, \texttt{def:consistency} & prove now \\
\texttt{lem:consistency-consequences} & \texttt{MIPStarRE/Quantum/Measurement.lean} & \texttt{def:consistency} & prove now \\
\texttt{lem:closeness-to-consistency} & \texttt{MIPStarRE/Quantum/Measurement.lean} & \texttt{def:consistency} & prove now \\
\texttt{lem:closeness-to-close-ips} & \texttt{MIPStarRE/Quantum/Measurement.lean} & \texttt{def:consistency} & prove now \\
\texttt{lem:add-a-proj} & \texttt{MIPStarRE/Quantum/Measurement.lean} & \texttt{def:consistency} & prove now \\
\texttt{lem:transfer-cons} & \texttt{MIPStarRE/Quantum/Measurement.lean} & \texttt{lem:closeness-to-close-ips} & prove now \\
\texttt{lem:cons-sub-meas} & \texttt{MIPStarRE/Quantum/Measurement.lean} & \texttt{def:consistency} & prove now \\
\texttt{lem:switcheroo} & \texttt{MIPStarRE/Quantum/Measurement.lean} & \texttt{lem:add-a-proj} & prove now \\
\texttt{def:code} & \texttt{MIPStarRE/Codes/LinearCode.lean} & --- & prove now \\
\texttt{prop:distance0} & \texttt{MIPStarRE/Codes/LinearCode.lean} & \texttt{def:code} & prove now \\
\texttt{def:interpolable} & \texttt{MIPStarRE/Codes/LinearCode.lean} & \texttt{prop:distance0} & prove now \\
\texttt{def:axis-line} & \texttt{MIPStarRE/Codes/LinearCode.lean} & --- & prove now \\
\texttt{def:tensor-code} & \texttt{MIPStarRE/Codes/LinearCode.lean} & \texttt{def:axis-line}, \texttt{def:code} & prove now \\
\texttt{prop:distance} & \texttt{MIPStarRE/Codes/LinearCode.lean} & \texttt{def:tensor-code}, \texttt{prop:distance0} & prove now \\
\texttt{prop:tuple-to-code-correspondence} & \texttt{MIPStarRE/Codes/LinearCode.lean} & \texttt{def:tensor-code}, \texttt{prop:distance0} & prove now \\
\texttt{prop:interpolate-tuple} & \texttt{MIPStarRE/Codes/LinearCode.lean} & \texttt{def:interpolable}, \texttt{prop:tuple-to-code-correspondence} & prove now \\
\texttt{def:tensor-code-test} & \texttt{MIPStarRE/Games/TensorCodeTest.lean} & \texttt{def:axis-line}, \texttt{def:subcube}, \texttt{def:tensor-code} & prove now \\
\texttt{def:tracial-strat} & \texttt{MIPStarRE/Games/TensorCodeTest.lean} & \texttt{def:tensor-code-test} & prove now \\
\texttt{def:tracial-good} & \texttt{MIPStarRE/Games/TensorCodeTest.lean} & \texttt{def:tracial-strat}, \texttt{def:consistency} & prove now \\
\texttt{prop:test-goodness} & \texttt{MIPStarRE/Games/TensorCodeTest.lean} & \texttt{def:tensor-code-test}, \texttt{def:tracial-good} & prove now \\
\texttt{prop:laplacian-edge-form} & \texttt{MIPStarRE/Quantum/Measurement.lean} & --- & prove now \\
\texttt{lem:local-to-global-expander} & \texttt{MIPStarRE/Quantum/Measurement.lean} & \texttt{prop:laplacian-edge-form} & prove now \\
\bottomrule
\end{longtable}
\normalsize

\subsection{Section~4--6 skeleton nodes}

\small
\begin{longtable}{L{0.18\textwidth}L{0.25\textwidth}L{0.32\textwidth}L{0.13\textwidth}}
\toprule
Node & Intended Lean home & Direct dependencies & Initial policy \\
\midrule
\endfirsthead
\toprule
Node & Intended Lean home & Direct dependencies & Initial policy \\
\midrule
\endhead
\texttt{lem:projectivization} & \texttt{MIPStarRE/Paper2111/Skeleton.lean} & \texttt{def:submeasurement}, \texttt{def:fd-trace} & placeholder \\
\texttt{lem:duality} & \texttt{MIPStarRE/Paper2111/Skeleton.lean} & \texttt{def:fd-trace} & placeholder \\
\texttt{lem:local-variance} & \texttt{MIPStarRE/Paper2111/Skeleton.lean} & \texttt{def:tracial-good}, \texttt{lem:consistency-consequences}, \texttt{prop:distance} & prove next \\
\texttt{lem:variance} & \texttt{MIPStarRE/Paper2111/Skeleton.lean} & \texttt{lem:local-variance}, \texttt{lem:local-to-global-expander} & prove next \\
\texttt{lem:self-improvement} & \texttt{MIPStarRE/Paper2111/Skeleton.lean} & \texttt{lem:projectivization}, \texttt{lem:duality}, \texttt{lem:variance} & placeholder \\
\texttt{lem:a-comm} & \texttt{MIPStarRE/Games/TensorCodeTest.lean} & \texttt{def:tracial-good}, \texttt{lem:add-a-proj} & prove next \\
\texttt{lem:g-comm-after-eval} & \texttt{MIPStarRE/Paper2111/Skeleton.lean} & \texttt{lem:a-comm}, \texttt{lem:cons-sub-meas} & placeholder \\
\texttt{lem:g-comm} & \texttt{MIPStarRE/Paper2111/Skeleton.lean} & \texttt{lem:g-comm-after-eval}, \texttt{prop:distance} & placeholder \\
\texttt{cor:G-hat-facts} & \texttt{MIPStarRE/Paper2111/Skeleton.lean} & \texttt{lem:g-comm} & placeholder \\
\texttt{lem:h-b-consistency} & \texttt{MIPStarRE/Paper2111/Skeleton.lean} & \texttt{cor:G-hat-facts}, \texttt{prop:interpolate-tuple} & placeholder \\
\texttt{cor:ha-cons} & \texttt{MIPStarRE/Paper2111/Skeleton.lean} & \texttt{lem:h-b-consistency}, \texttt{lem:data-processing}, \texttt{lem:transfer-cons} & placeholder \\
\texttt{lem:over-all-outcomes} & \texttt{MIPStarRE/Paper2111/Skeleton.lean} & \texttt{prop:interpolate-tuple}, \texttt{prop:distance} & placeholder \\
\texttt{lem:from-H-to-G} & \texttt{MIPStarRE/Paper2111/Skeleton.lean} & \texttt{cor:G-hat-facts} & placeholder \\
\texttt{lem:chernoff-bernoulli-matrix} & \texttt{MIPStarRE/Paper2111/Skeleton.lean} & \texttt{def:fd-trace} & placeholder \\
\texttt{cor:ld-pasting-N-completeness} & \texttt{MIPStarRE/Paper2111/Skeleton.lean} & \texttt{lem:over-all-outcomes}, \texttt{lem:from-H-to-G}, \texttt{lem:chernoff-bernoulli-matrix} & placeholder \\
\texttt{lem:pasting} & \texttt{MIPStarRE/Paper2111/Skeleton.lean} & \texttt{cor:ha-cons}, \texttt{cor:ld-pasting-N-completeness} & placeholder \\
\texttt{lem:induction} & \texttt{MIPStarRE/Paper2111/Skeleton.lean} & \texttt{lem:self-improvement}, \texttt{lem:pasting} & placeholder \\
\texttt{thm:main} & \texttt{MIPStarRE/Paper2111/Skeleton.lean} & \texttt{lem:induction}, \texttt{lem:self-improvement} & placeholder \\
\texttt{def:good-2} & \texttt{MIPStarRE/Paper2111/Skeleton.lean} & \texttt{def:tensor-code-test}, \texttt{def:tracial-good} & placeholder \\
\texttt{thm:main-bipartite} & \texttt{MIPStarRE/Paper2111/Skeleton.lean} & \texttt{thm:main}, \texttt{def:good-2} & placeholder \\
\bottomrule
\end{longtable}
\normalsize

\subsection{First formalization tasks}

The first fifteen Lean tasks, in dependency order, should be the following.

\small
\begin{longtable}{L{0.05\textwidth}L{0.49\textwidth}L{0.23\textwidth}L{0.17\textwidth}}
\toprule
\# & Task & Target file & Main unlocks \\
\midrule
\endfirsthead
\toprule
\# & Task & Target file & Main unlocks \\
\midrule
\endhead
1 & Define the finite-dimensional normalized trace and the trace seminorm on matrices. & \texttt{Quantum/Measurement.lean} & all operator estimates \\
2 & Define measurements, submeasurements, projective measurements, and processed measurements. & \texttt{Quantum/Measurement.lean} & all consistency lemmas \\
3 & Define consistency and closeness for question-indexed operator families. & \texttt{Quantum/Measurement.lean} & Section~2 measurement API \\
4 & Prove data processing for consistency and consistency-to-closeness. & \texttt{Quantum/Measurement.lean} & test bookkeeping, local variance \\
5 & Prove closeness-to-consistency, overlap stability, multiplication stability, transfer of consistency, and the submeasurement control lemma. & \texttt{Quantum/Measurement.lean} & nearly every later analytic proof \\
6 & Prove the switcheroo lemma for products of approximately commuting operators. & \texttt{Quantum/Measurement.lean} & Section~6 commutation chain \\
7 & Define linear codes, their distance parameter, and the uniqueness-from-distance lemma. & \texttt{Codes/LinearCode.lean} & interpolation theory \\
8 & Define interpolable codes and the interpolation map. & \texttt{Codes/LinearCode.lean} & tuple interpolation \\
9 & Define axis-parallel lines and tensor codes. & \texttt{Codes/LinearCode.lean} & all paper-specific code objects \\
10 & Prove the tensor-code distance bound, with a correct induction or generator-matrix proof. & \texttt{Codes/LinearCode.lean} & local variance, commutation averaging \\
11 & Prove tuple-to-code correspondence for $\code^{\otimes(m+1)}$. & \texttt{Codes/LinearCode.lean} & pasting and interpolation \\
12 & Prove interpolation of tuples for interpolable base codes. & \texttt{Codes/LinearCode.lean} & Section~6 pasting skeleton \\
13 & Define subcubes, the augmented tensor code test, finite-dimensional tracial strategies, and $(\eps,\delta)$-good strategies. & \texttt{Games/TensorCodeTest.lean} & all paper statements \\
14 & Prove the winning-probability versus goodness bookkeeping lemma, then add named placeholder declarations for \texttt{lem:self-improvement}, \texttt{lem:pasting}, \texttt{lem:induction}, and \texttt{thm:main}. & \texttt{Games/TensorCodeTest.lean}, \texttt{Paper2111/Skeleton.lean} & executable paper skeleton \\
15 & Formalize the appendix Laplacian identity and local-to-global expander inequality. & \texttt{Quantum/Measurement.lean} & \texttt{lem:variance} and the first post-pilot theorem chain \\
\bottomrule
\end{longtable}
\normalsize

The first non-placeholder theorem after the pilot should be \Cref{lem:a-comm}. It is the shortest genuine bridge from the Section~3 test interface to the Section~6 pasting analysis. After that, \Cref{lem:local-variance} and \Cref{lem:variance} are the right next targets, because they exercise the appendix lemma in the exact way later used by self-improvement.

The initial placeholder policy is therefore as follows. Every node through \Cref{lem:local-to-global-expander,prop:test-goodness} should be proved immediately. The next-wave targets are \Cref{lem:a-comm,lem:local-variance,lem:variance}. The hard statements \Cref{lem:projectivization,lem:duality,lem:self-improvement,lem:g-comm-after-eval,lem:g-comm,lem:h-b-consistency,lem:over-all-outcomes,lem:from-H-to-G,lem:chernoff-bernoulli-matrix,lem:pasting,lem:induction,thm:main,thm:main-bipartite} should initially remain named placeholders or axioms in \texttt{MIPStarRE/Paper2111/Skeleton.lean}.

\end{document}